We present Gemini, a family of highly capable multimodal models developed at Google. We trained Gemini models jointly across image, audio, video, and text data for the purpose of building a model with both strong generalist capabilities across modalities alongside cutting-edge understanding and reasoning performance in each respective domain. 

Gemini 1.0, our first version, comes in three sizes: Ultra for highly-complex tasks, Pro for enhanced performance and deployability at scale, and Nano for on-device applications. Each size is specifically tailored to address different computational limitations and application requirements. 

After large-scale pre-training, we post-train our models to improve overall quality, enhance target capabilities, and ensure alignment and safety criteria are met. Due to the varied requirements of our downstream applications, we have produced two post-trained Gemini model family variants. Chat-focused variants, referred to as Gemini Apps models, are optimized for \href{https://gemini.google.com}{Gemini and Gemini Advanced}, our conversational AI service formerly known as Bard. Developer-focused variants, referred to as Gemini API models, are optimized for a range of products and are accessible through \href{https://makersuite.google.com}{Google AI Studio} and \href{https://cloud.google.com/vertex-ai}{Cloud Vertex AI}.

We evaluate the performance of pre- and post-trained Gemini models on a comprehensive suite of internal and external benchmarks covering a wide range of language, coding, reasoning, and multimodal tasks.

The Gemini family advances state-of-the-art in large-scale language modeling \citep{gpt2, gpt3, gopher, palm, chinchilla, palm2, gpt4}, image understanding \citep{vit, pali, coca, flamingo, gato, gpt4v}, audio processing \citep{whisper, usm}, and video understanding \citep{flamingo, palix}. It also builds on the work on sequence models \cite{sutskever2014sequence}, a long history of work in deep learning based on neural networks~\citep{lecun2015deep}, and machine learning distributed systems \cite{distbelief,jax,pathways} that enable large-scale training. 

Our most capable model, Gemini Ultra, achieves new state-of-the-art results in 30 of 32 benchmarks we report on, including 10 of 12 popular text and reasoning benchmarks, 9 of 9 image understanding benchmarks, 6 of 6 video understanding benchmarks, and 5 of 5 speech recognition and speech translation benchmarks. Gemini Ultra is the first model to achieve human-expert performance on MMLU \citep{mmlu} –- a prominent benchmark testing knowledge and reasoning via a suite of exams –- with a score above 90\%.
Beyond text, Gemini Ultra makes notable advances on challenging multimodal reasoning tasks. For example, on the recent MMMU benchmark~\cite{mmmu}, that comprises questions about images on multi-discipline tasks requiring college-level subject knowledge and deliberate reasoning, Gemini Ultra achieves a new state-of-the-art score of 62.4\%, outperforming the previous best model by more than 5 percentage points. It provides a uniform performance lift for video question answering and audio understanding benchmarks.

Qualitative evaluation showcases impressive crossmodal reasoning capabilities, enabling the model to understand and reason across an input sequence of audio, images, and text natively (see Figure~\ref{fig:inverse_graphics_3} and Table~\ref{tab:modality-combo}). Consider the educational setting depicted in Figure~\ref{fig:homeworkexample} as an example. A teacher has drawn a physics problem of a skier going down a slope, and a student has worked through a solution to it. Using Gemini models' multimodal reasoning capabilities, the model is able to understand the messy handwriting, correctly understand the problem formulation, convert both the problem and solution to mathematical typesetting, identify the specific step of reasoning where the student went wrong in solving the problem, and then give a worked through correct solution to the problem. This opens up exciting educational possibilities, and we believe the new multimodal and reasoning capabilities of Gemini models have dramatic applications across many fields.


\begin{figure*}[h!]
\centering
\includegraphics[width=\textwidth,keepaspectratio]{figs/homework_physics_problem.jpg}
\caption{Verifying a student's solution to a physics problem. The model is able to correctly recognize all of the handwritten content and verify the reasoning. On top of understanding the text in the image, it needs to understand the problem setup and correctly follow instructions to generate \LaTeX.}
\label{fig:homeworkexample}
\end{figure*}

The reasoning capabilities of large language models show promise toward building generalist agents that can tackle more complex multi-step problems. The AlphaCode team built AlphaCode 2~\citep{gdm2023alphacode2}, a new Gemini-model-powered agent, that combines Gemini models' reasoning capabilities with search and tool-use to excel at solving competitive programming problems. 
AlphaCode 2 ranks within the top 15\% of entrants on the Codeforces competitive programming platform, a large improvement over its state-of-the-art predecessor in the top 50\% \citep{alphacode}. 

In tandem, we advance the frontier of efficiency with Gemini Nano, a series of small models targeting on-device deployment. These models excel in on-device tasks, such as summarization, reading comprehension, text completion tasks, and exhibit impressive capabilities in reasoning, STEM, coding, multimodal, and multilingual tasks relative to their sizes. 

In the following sections, we first provide an overview of the model architecture, training infrastructure, and pre-training dataset. We then present detailed \textit{evaluations} of the pre- and post-trained Gemini model family, covering well-studied benchmarks across text, code, image, audio and video --- which include both English performance and multilingual capabilities.  Next we discuss our approach to post-training, highlight common and distinct aspects of the Gemini Apps and Gemini API model variants, and benchmark their performance on key capabilities. \textit{Responsible deployment} is critical: we explain our process for impact assessments, developing model policies, evaluations, and mitigations of harm before deployment decisions. Finally, we discuss the broader implications of Gemini models, their limitations alongside their potential applications --- paving the way for a new era of research and innovation in AI. 
